\section{Gaussian Space --- Hermite --- Foundations}
%%% generated by the blueprint pipeline from TCSlib.BooleanAnalysis.GaussianSpace.Hermite.Basic
\subsection*{Gaussian Space --- Hermite --- Foundations: Introduction}

Declarations of the Lean module \texttt{TCSlib.BooleanAnalysis.GaussianSpace.Hermite.Basic}.

\subsection*{Gaussian Space --- Hermite --- Foundations: Results}

\begin{theorem}[Zeroth normalized Hermite polynomial]
\lean{BooleanAnalysis.GaussianSpace.normalizedHermitePolynomial_zero}
\label{BooleanAnalysis.GaussianSpace.normalizedHermitePolynomial_zero}
\leanok
\uses{BooleanAnalysis.GaussianSpace.normalizedHermitePolynomial}
\difficulty{1}
The normalized Hermite polynomial of degree zero is the constant polynomial $1$.
\end{theorem}

\begin{theorem}[First normalized Hermite polynomial]
\lean{BooleanAnalysis.GaussianSpace.normalizedHermitePolynomial_one}
\label{BooleanAnalysis.GaussianSpace.normalizedHermitePolynomial_one}
\leanok
\uses{BooleanAnalysis.GaussianSpace.normalizedHermitePolynomial}
\difficulty{1}
The normalized Hermite polynomial of degree one is the identity polynomial $X$.
\end{theorem}

\begin{theorem}[Product formula for Hermite functions]
\lean{BooleanAnalysis.GaussianSpace.hermiteFunction_eq_prod}
\label{BooleanAnalysis.GaussianSpace.hermiteFunction_eq_prod}
\leanok
\uses{BooleanAnalysis.GaussianSpace.hermiteFunction, BooleanAnalysis.GaussianSpace.normalizedHermitePolynomial}
\difficulty{3}
For every natural number $n$, multi-index $\alpha:\{1,\ldots,n\}\to\mathbb N$, and point
$x\in\mathbb R^n$, the multivariate Hermite function is the product of its one-variable
normalized Hermite polynomials: $H_\alpha(x)=\prod_{i=1}^n h_{\alpha_i}(x_i)$.
\end{theorem}

\begin{theorem}[Hermite function at the zero index]
\lean{BooleanAnalysis.GaussianSpace.hermiteFunction_zero}
\label{BooleanAnalysis.GaussianSpace.hermiteFunction_zero}
\leanok
\uses{BooleanAnalysis.GaussianSpace.hermiteFunction, BooleanAnalysis.GaussianSpace.hermiteFunction_eq_prod}
\difficulty{1}
For every natural number $n$ and every $x\in\mathbb R^n$, the Hermite function indexed
by the zero multi-index has value $1$.
\end{theorem}

\begin{theorem}[Hermite function for an indicator index]
\lean{BooleanAnalysis.GaussianSpace.hermiteFunction_indicator}
\label{BooleanAnalysis.GaussianSpace.hermiteFunction_indicator}
\leanok
\uses{BooleanAnalysis.GaussianSpace.hermiteFunction, BooleanAnalysis.GaussianSpace.hermiteFunction_eq_prod, BooleanAnalysis.GaussianSpace.normalizedHermitePolynomial_one, BooleanAnalysis.GaussianSpace.normalizedHermitePolynomial_zero}
\difficulty{1}
For every natural number $n$, finite set $S\subseteq\{1,\ldots,n\}$, and point
$x\in\mathbb R^n$, the Hermite function indexed by the multi-index that is $1$ on $S$
and $0$ elsewhere equals $\prod_{i\in S}x_i$.
\end{theorem}

\begin{lemma}[Characterization of an Indicator Multi-Index]
\lean{BooleanAnalysis.GaussianSpace.indicatorMultiIndex_eq_iff}
\label{BooleanAnalysis.GaussianSpace.indicatorMultiIndex_eq_iff}
\leanok
\difficulty{4}
Let $n$ be a natural number, let $S\subseteq\{0,\ldots,n-1\}$, and let
$\alpha\colon\{0,\ldots,n-1\}\to\mathbb N$. The function that takes value $1$ on $S$ and
$0$ off $S$ equals $\alpha$ if and only if every value of $\alpha$ is at most $1$ and
$S=\{i:\alpha(i)=1\}$.
\end{lemma}

\begin{lemma}[Derivative of a Gaussian Hermite polynomial]
\lean{BooleanAnalysis.GaussianSpace.derivative_hermite_succ}
\label{BooleanAnalysis.GaussianSpace.derivative_hermite_succ}
\leanok
\difficulty{4}
For every natural number $k$, the derivative of the Gaussian Hermite polynomial of
degree $k+1$ is $(k+1)$ times the Gaussian Hermite polynomial of degree $k$.
\end{lemma}
